\centerline{\bf \Huge Первый разнобой 2025}

\begin{flushright}
        \scriptsize{Школа — это лишь тесный аквариум,\\ а в мире есть океан и другие аквариумы. Выбери своё. \\Например, техлицей...\\Наги Ёсино. Магическая битва.}
\end{flushright}

\problem{1}
Положительные числа $a, b, c$ известно, что

$$
a^2+b^2+c^2=2 a b+2 a c+2 b c
$$


Докажите, что среди чисел $\sqrt{a}, \sqrt{b}, \sqrt{c}$ одно равно сумме двух других.


\problem{2}
На плоскости дана прямая и 101 точка, причем расстояние от любой точки до прямой не превосходит 1 , а расстояние между любыми двумя точками не менее 2. Докажите, что какие-то две точки находятся на расстоянии более 85.

\problem{3}
Дан квадратный трехчлен $P(x)$. Докажите, что существуют попарно различные числа $a, b$ и $c$ такие, что выполняются равенства

$$
P(b+c)=P(a), P(c+a)=P(b), P(a+b)=P(c) .
$$

\problem{4}
На стороне $C D$ квадрата $A B C D$ выбрана точка $P$. Точка $Q$ на стороне $A D$ такова, что $B Q$ - биссектриса угла $P B A$. Из точки Р на $B Q$ опущен перпендикуляр $P H$. Докажите, что $B Q=2 P H$.

\problem{5}
Обозначим через $d(k)$ количество натуральных делителей натурального числа $k$. Докажите, что для каждого натурального $n$

$$
d(1)+d(3)+d(5)+\ldots+d(2 n-1)<d(2)+d(4)+d(6)+\ldots+d(2 n) .
$$

\problem{6}
Пусть $C E$ - биссектриса в остроугольном треугольнике $A B C$. На внешней биссектрисе угла $A C B$ отмечена точка $D$, а на стороне $B C$ - точка $F$, причём $\angle B A D=90^{\circ}=\angle D E F$. Докажите, что центр окружности, описанной около треугольника $CEF$,
лежит на прямой $BD$. 
